\documentclass[a4paper]{article}

\usepackage[english]{babel}
\usepackage[utf8]{inputenc}
\usepackage[hidelinks]{hyperref}

\usepackage{graphicx}
\graphicspath{ {images/} }

\usepackage{setspace}
\usepackage{pdfpages}
\usepackage{xspace} %Removes space after commands
\usepackage{fancyhdr}
\usepackage{lastpage}
\usepackage{verbatim}
\usepackage{url}

%%% BibTeX %%%
\usepackage{cite}
% JS function to get BibTeX entry from WebAssembly website
\begin{comment}
function BibTeX() {
	let subject = document.title.replace(/ - WebAssembly/, '').replace(' ', '-').toLowerCase();
    return `@misc{website:wasm-${subject}
	author = "The WebAssembly working group",
	title = {{${document.title}}},
	url = "${document.URL}",
	note = {Online; accessed 3 May 2017},
}`;
}
\end{comment}

% source: https://tex.stackexchange.com/a/35044
\makeatletter
\newcommand\footnoteref[1]{\protected@xdef\@thefnmark{\ref{#1}}\@footnotemark}
\makeatother

% Code highlighting
\usepackage{minted}
% Highlight theme
\usemintedstyle{trac}

\pagestyle{fancy}
\fancyhf{}

\lfoot{Page \thepage \hspace{1pt} of \pageref{LastPage}}

\usepackage[margin=0.8in]{geometry}

\title{Bachelor Project Report: Compiling MicroC to WebAssembly}

\begin{document}
%\pagenumbering{gobble} %Remove page numbers
\begin{titlepage}
	\centering
	{\scshape\LARGE IT University of Copenhagen \par}
	\vspace{1cm}
	{\scshape\Large Bachelor Project Report\par}
	\vspace{1.5cm}
	{\huge\bfseries Compiling MicroC to WebAssembly \par}
	\vspace{2cm}
	{\includegraphics{web-assembly-logo-black-512px} \par}
	\vspace{2cm}
	{\Large\itshape Andreas Bjørn Hassing Nielsen\par}
	abhn@itu.dk\\
	\vspace{2cm}
	{\Large Abstract\par}
	%TODO: Write an awesome abstract.
	{\bfseries This is where the abstract will go.}
	\vfill
% Bottom of the page
	{\large \today\par}
\end{titlepage}

\newpage

\tableofcontents
\newpage


\section{Introduction}
The purpose of this project is to build a compiler, also called a translator, from MicroC to WebAssembly, using FsLexYacc\footnote{\label{footnote:fslexyacc-url}http://fsprojects.github.io/FsLexYacc/} and F\#\footnote{http://fsharp.org/}. The learning goals are to develop an understanding of WebAssembly and improve my knowledge of compiler design.

MicroC is a C-like language described by Peter Sestoft in the book Programming Language Concepts~\cite{PLC}. The original lexer and parser code can be found at his website\footnote{http://www.itu.dk/people/sestoft/plc/}.


\subsection{WebAssembly}
Semantics~\cite{website:wasm-semantics}.

\subsection{Compilers}
What is a compiler? According to the definition laid out by Mogensen~\cite{BCD}, a compiler is a program that takes another program, written in some higher-level language, C\# for instance, and compiles it to a lower-level representation, Common Intermediate Language code - in the case of C\#.

Languages running on the .NET CLR, F\# and C\# to name two, compiles to bytecode known as Common Intermediate Language (CIL), which is then compiled to the specific instruction set on the machine running the program just-in-time (JIT), whereas languages such as C and Rust are compiled ahead-of-time directly to the instruction set of the compiling machine\footnote{Unless the compiler is configured to compile to another OS and/or processor architecture.}.

For this report, only AOT compilation will be discussed.

\section{Problem Definition}
WebAssembly (WASM) is a new up-and-coming binary code format for the user-facing web, designed to work alongside JavaScript as a more portable, size- and load-time-efficient, ahead-of-time-compiled alternative. 

The binary WASM format is not bound to be emitted by JavaScript (asm.js) only, potentially bringing your favourite language to the browser front-end. If a language can compile to a intermediate representation supported by WASM (LLVM IR, for instance), it can be compiled to the binary WASM format. 

The purpose of this project is to shed some light on the unfinished WebAssembly specification. To assist in meeting this goal, a compiler will be designed, that compiles from a simple programming language, MicroC, to a WASM format that can be run in a WASM-enabled browser today (currently requiring Firefox Nightly or Chrome Canary with WASM flags set to enabled\footnote{At the time of writing (May, 2017), setting flags or using a nightly browser version is no longer required. Major browsers now support WASM out of the box: https://lists.w3.org/Archives/Public/public-webassembly/2017Feb/0002.html}).

\subsection{Method}
What follows is the planned activities and sources of information.
\begin{itemize}
	\item Translation of MicroC to the binary WebAssembly format, using FsLexYacc\footnoteref{footnote:fslexyacc-url}, as used in the Programs as Data course (fall 2016).
	\item Course books from Programs as Data (fall 2016) will be used as knowledge base for MicroC to WebAssembly translation:
		\begin{itemize}
			\item Peter Sestoft: Programming Language Concepts. Springer 2012.
			\item Torben Mogensen: Basics of Compiler Design. DIKU 2010. Chapters 2 and 3.
		\end{itemize}
	\item Knowledge of WebAssembly will be gained through the official website (http://webassembly.org/) and links branching out from it.
	\item Creation of prototype: web interface that lets users type MicroC in a window, and see it generate WebAssembly in some human-readable format, either s-expression or linear bytecode. The prototype will also give users the ability to run their code as WebAssembly, given that the browser in use supports it.
	\item Reflection over, either the security behind WebAssembly or the portability of it between browsers and platforms.
\end{itemize}

\noindent If time allows:
\begin{itemize}
	\item Extensions to the MicroC language will be implemented.
	\item Simple optimizations to the generated WebAssembly code.
\end{itemize}

\section{Problem Analysis}

\section{User Guide and Examples}

\section{Technical Description}

\section{Testing and Validation}

\section{Peripheral Tooling}

\section{Extensions}

\section{Conclusion}

\section{References}
\begingroup
\renewcommand{\section}[2]{}%
\bibliographystyle{plain}
\bibliography{Bibliography}{}
\endgroup

\section{Appendix}

\end{document}
